\documentclass[12pt]{article}

% ── packages ──────────────────────────────────────────────────────────────────
\usepackage{amsmath,amssymb,amsthm}
\usepackage{geometry}
\usepackage{hyperref}
\usepackage{booktabs}
\usepackage{enumitem}
\usepackage{parskip}

\geometry{letterpaper, margin=1in}

% ── theorem environments ───────────────────────────────────────────────────────
\newtheorem{theorem}{Theorem}[section]
\newtheorem{definition}[theorem]{Definition}
\newtheorem{postulate}[theorem]{Postulate}
\newtheorem{corollary}[theorem]{Corollary}
\newtheorem{remark}[theorem]{Remark}

% ── convenience macros ────────────────────────────────────────────────────────
\newcommand{\lP}{l_P}
\newcommand{\EP}{E_P}
\newcommand{\tP}{t_P}
\newcommand{\SSV}{\text{SSV}}
\newcommand{\PSR}{\text{PSR}}
\newcommand{\DP}{\text{DP}}
\newcommand{\CP}{\text{CP}}
\newcommand{\GP}{\text{GP}}
\newcommand{\alphageom}{\alpha_{\rm geom}}
\newcommand{\SSVcrit}{\SSV_{\rm crit}}

% ── title block ───────────────────────────────────────────────────────────────
\title{\textbf{Microscopic Origin of the Dipole Sea Stiffness \boldmath$C$\\
in Conscious Point Physics}\\[6pt]
{\large Companion Paper to ``Mechanistic Derivation of Relativistic Effects\\
via Space Stress Vector (SSV) in the Dipole Sea'' (Version~16)}}

\author{Thomas Lee Abshier, ND\\
Hyperphysics Institute\\
\texttt{https://hyperphysics.com}\\
\texttt{drthomas007@protonmail.com}}

\date{13 March 2026 --- Version~2.1}

\begin{document}
\maketitle

% ── keywords ──────────────────────────────────────────────────────────────────
\noindent\textbf{Keywords:} Conscious Point Physics, dipole sea stiffness,
Voronoi second-moment integral, 600-cell geometry, H4 symmetry,
inverse-square force law, fractional quark charges, cage-occupation effect,
Planck-scale Coulomb interaction.

% ── abstract ──────────────────────────────────────────────────────────────────
\begin{abstract}
The elastic stiffness $C$ of the Dipole Sea governs the linear Hooke-like
response of Space Stress Vector (SSV) and appears in the main paper's central
formula $\PSR_{\rm eff} = \lP/(1 + k\cdot\Delta\SSV)$ as $k = 1/C\cdot\EP/\lP^3$.
We establish the microscopic origin of $C$ in three steps.  First, all four
Dipole Particle (DP) types present in the equilibrium sea carry the same
elementary charge $\pm 1$ with identical Coulombic coupling strength; their
equal abundances and the quasi-static limit of Zitterbewegung (ZBW)
oscillations together reduce the multi-component medium to a single effective
stiffness.  Second, that stiffness is the second-moment integral of the
600-cell Voronoi face-area distribution over all SSV directions, computed
exactly from the H4 golden-ratio coordinates in Version~16 Appendix~A.5:
\[
  C \;=\; \alphageom \cdot \SSVcrit,\qquad
  \SSVcrit = \frac{\EP}{\lP^3},\qquad
  \alphageom \;=\; \frac{3(11+5\sqrt{5})\sqrt{5+\sqrt{5}}}{320}
  \;\approx\; 0.5594.
\]
Third, the identical stiffness $C$ and the 12-neighbor shell-broadcast
geometry of SSV quanta together reproduce the universal inverse-square force
law and, via cage-occupation projection factors fixed by H4 representation
theory, predict the fractional effective charges of quarks relative to
leptons --- a result deferred in detail to a dedicated companion paper.
No new postulates are introduced beyond those of Version~16.
\end{abstract}

% ══════════════════════════════════════════════════════════════════════════════
\section{Introduction}
\label{sec:intro}
% ══════════════════════════════════════════════════════════════════════════════

The central formula of Version~16 is
\begin{equation}
  \PSR_{\rm eff} \;=\; \frac{\lP}{1 + k\cdot\Delta\SSV},
  \qquad k = \frac{\lP^3}{\EP} \approx 2.16\times10^{-114}\ \text{m}^3/\text{J}.
  \label{eq:PSR}
\end{equation}
The coupling constant $k$ encodes the elastic stiffness $C = \EP/(\alphageom\lP^3)$
of the Dipole Sea against SSV perturbation.  In Version~16 this value is
fixed by dimensional necessity alone: $k$ is the unique combination of Planck
constants with units m$^3$/J.  That is a necessary condition but not a
sufficient one; a referee may reasonably ask whether the microscopic
Planck-scale interactions of the DP sea actually produce this value, or whether
a different stiffness is equally consistent with the CPP postulates.

This companion paper answers that question.  We show that:
\begin{enumerate}
  \item the four DP types present in the equilibrium sea all carry the same
    elementary Coulombic coupling (Section~\ref{sec:dp-sea}),
  \item their combined stiffness is the H4 Voronoi second-moment integral
    already computed in Appendix~A.5 of Version~16
    (Section~\ref{sec:voronoi}), and
  \item the same stiffness and broadcast geometry reproduce the inverse-square
    force law and suggest the origin of fractional quark charges
    (Section~\ref{sec:implications}).
\end{enumerate}
No new postulates are required.

% ══════════════════════════════════════════════════════════════════════════════
\section{The Equilibrium Dipole Sea}
\label{sec:dp-sea}
% ══════════════════════════════════════════════════════════════════════════════

\subsection{Four DP types with equal abundances}

A Dipole Particle is a bound pair of two Conscious Points of opposite polarity.
The four distinct DP types are:
\begin{center}
\begin{tabular}{@{}clll@{}}
\toprule
Type & Constituents & Coupling & Abundance \\
\midrule
1 & $\text{eCP}^+/\text{eCP}^-$ & electron-type & $n_0$ \\
2 & $\text{qCP}^+/\text{qCP}^-$ & quark-type    & $n_0$ \\
3 & $\text{qCP}^+/\text{eCP}^-$ & hybrid (A)    & $n_0$ \\
4 & $\text{eCP}^+/\text{qCP}^-$ & hybrid (B)    & $n_0$ \\
\bottomrule
\end{tabular}
\end{center}
The equal abundance $n_0$ per type follows from the symmetry of the
equilibrium state: the DP sea was formed from an equal mixture of all CP
types, and maximum-entropy equilibrium distributes pairings uniformly across
all four combinations.

\subsection{Identical Coulombic coupling for all CP types}

Every Conscious Point carries the same elementary charge state $\pm 1$
(Version~16, \S 2).  There is no distinction at the CP level between
``electron-type'' and ``quark-type'' in terms of fundamental charge magnitude.
The apparent difference between quarks and leptons arises at the \emph{composite}
level from cage-occupation geometry (Section~\ref{sec:implications}), not from
any difference in CP charge strength.

Therefore, to leading order, all four DP types respond to an external SSV
perturbation with the \emph{same} Coulombic coupling constant.  The
multi-component stiffness sum
\begin{equation}
  C_{\rm total} \;=\; \sum_{i=1}^{4} n_i\, C_i
  \;=\; 4\, n_0\, C_1
  \label{eq:Ctotal}
\end{equation}
reduces to a single-component problem with total density $4n_0$ and the
uniform single-pair stiffness $C_1$.

\subsection{Quasi-static limit of ZBW oscillations}

Each DP undergoes Zitterbewegung (ZBW) oscillation at frequency
$\nu_{\rm ZBW} = (2\tP)^{-1}$, alternating between attraction to its
pair-mate and attraction to the surrounding DP sea.  The SSV response relevant
to special-relativistic effects operates at macroscopic timescales
$\tau \gg \tP$, i.e.\ at frequencies $\nu \ll \nu_{\rm ZBW}$.  In this
quasi-static limit the time-average of the ZBW correction factor is unity:
\begin{equation}
  \langle f_{\rm ZBW}\rangle \;=\;
  \frac{1}{T}\int_0^T [1 + \beta\cos(2\pi\nu_{\rm ZBW} t)]\,dt
  \;=\; 1,
  \label{eq:ZBW}
\end{equation}
so ZBW oscillations do not affect the stiffness at SR timescales.  Their role
in quantum-scale phenomena (ZBW as the origin of $\hbar$) is deferred to the
Born-rule companion paper.

\subsection{Maximum-entropy equilibrium}

The equilibrium DP sea has maximum entropy: DP orientations are uniformly
random, there is no net SSV in any direction, and no persistent local
ordering.  This isotropy is a prerequisite for Lorentz covariance and is
maintained dynamically by the CP Exclusion Rule (if two CPs occupy the same
Grid Point at the same $t_{\rm abs}$, both instantly displace to the PSR edge,
randomizing local configurations and preventing ordered clustering).

% ══════════════════════════════════════════════════════════════════════════════
\section{The Voronoi Second-Moment Integral}
\label{sec:voronoi}
% ══════════════════════════════════════════════════════════════════════════════

\subsection{Structure of the 600-cell Voronoi cell}

Each Grid Point in the 600-cell lattice occupies the center of a Voronoi
cell bounded by 12 pentagonal faces, whose outward normals point toward the
12 nearest-neighbor GPs (coordination number 12).  The Voronoi cell geometry
is fully determined by the golden-ratio H4 coordinates of the 600-cell
(Version~16, \S A.2).

\subsection{Definition of the stiffness integral}

When the DP sea at a GP is displaced by an infinitesimal SSV perturbation
$\delta\SSV$ in direction $\hat{\mathbf{r}}$, the restoring force per unit
volume is proportional to the second moment of the Voronoi face-area
distribution projected onto $\hat{\mathbf{r}}$:
\begin{equation}
  C \;=\; \frac{1}{V_0}\sum_{i=1}^{12} A_i
  \bigl\langle(\hat{n}_i\cdot\hat{\mathbf{r}})^2\bigr\rangle_{\hat{\mathbf{r}}},
  \label{eq:Cintegral}
\end{equation}
where $A_i$ is the area of the $i$-th pentagonal face, $\hat{n}_i$ is its
outward unit normal, and $V_0$ is the Voronoi cell volume.  The angular
average is taken uniformly over all SSV directions $\hat{\mathbf{r}}$.

\subsection{Exact evaluation from H4 coordinates}

The exact evaluation of~\eqref{eq:Cintegral} requires the full 4D golden-ratio
coordinate set of the 600-cell, as carried out in Version~16, Appendix~A.5.
The result is
\begin{equation}
  \boxed{C \;=\; \alphageom\cdot\SSVcrit,\qquad
  \SSVcrit = \frac{\EP}{\lP^3},\qquad
  \alphageom = \frac{3(11+5\sqrt{5})\sqrt{5+\sqrt{5}}}{320}
  \approx 0.5594.}
  \label{eq:C}
\end{equation}
This is an exact algebraic constant, fixed entirely by H4 symmetry.

\begin{remark}[Why no simple closed form suffices]
One might expect $C$ to take the simple form $m\bar{A}/V_0$ for some small
integer $m$, where $\bar{A}$ is the mean pentagonal face area.  Numerical
computation shows that neither $m=3$ (the 4D angular average $1/4$ applied to
12 faces) nor $m=4$ (the 3D angular average $1/3$ applied to 12 faces)
reproduces $\alphageom\cdot\EP/\lP^3$ when the Voronoi cell is treated as a
regular dodecahedron.  The exact value~\eqref{eq:C} requires the full H4
integral because the 600-cell Voronoi faces are \emph{not} the regular
pentagons of an ideal dodecahedron --- they are distorted by the 4D embedding
and the golden-ratio vertex positions.  Version~16 Appendix~A.5 contains the
complete algebraic derivation; this companion paper takes the result as given
and explores its physical consequences.
\end{remark}

\subsection{Interpretation: $C$ as the unique Planck-scale stiffness}

Equation~\eqref{eq:C} has a simple dimensional interpretation.
$\SSVcrit = \EP/\lP^3$ is the Planck energy density --- the maximum
SSV storable in one Planck volume without geometric collapse.
The factor $\alphageom \approx 0.5594 < 1$ is the geometric efficiency
of the 600-cell's icosahedral packing: the lattice can only exploit
$\approx 56\%$ of the theoretical maximum Planck stiffness because its
12-fold coordination and pentagonal faces do not tile 3D space perfectly.
No other regular 4-polytope produces this precise efficiency; the 600-cell
is selected by the same argument that fixes the dimensionality theorem in
Version~16 (Appendix~G).

% ══════════════════════════════════════════════════════════════════════════════
\section{Physical Implications}
\label{sec:implications}
% ══════════════════════════════════════════════════════════════════════════════

\subsection{Universal inverse-square force law}

With $C$ fixed by~\eqref{eq:C}, the SSV strain at distance $r$ from a
source CP is
\begin{equation}
  \varepsilon(r) \;=\; \frac{\Delta\SSV(r)}{C}
  \;=\; k\cdot\Delta\SSV(r).
  \label{eq:strain}
\end{equation}
The SSV field itself obeys the inverse-square law because each GP broadcasts
its fixed SSV quantum $Q$ to the shell of GPs at radius $r = \PSR$, one shell
per Absolute Moment tick (Absolute Moment companion paper, \S 5.2).
The broadcast follows the 12-neighbor edge selection already used in the
Absolute Moment companion to imprint left-handed chirality on CP transit
(Absolute Moment companion, \S 3, Eq.~(3)); the same 12-edge geometry that
drives chirality also drives field propagation.
The fixed quantum $Q$ is distributed over the $\sim 4\pi r^2$ GPs on the
shell, giving
\begin{equation}
  \Delta\SSV(r) \;\propto\; \frac{Q}{4\pi r^2}
  \;\propto\; \frac{1}{r^2}.
  \label{eq:invsq}
\end{equation}
The force between two source CPs therefore follows $F \propto q_1 q_2/r^2$
with a universal prefactor set entirely by $C$ and $Q$.  No separate
electromagnetic coupling constant is introduced; the Coulomb constant $k_e$
is determined by $C$, $Q$, and the Planck units already fixed in Version~16.

\subsection{Consistency with the PSR formula}

The strain~\eqref{eq:strain} directly supplies the denominator of the PSR
formula~\eqref{eq:PSR}.  Because $C$ (and therefore $k = 1/(C\lP^3/\EP)$)
is fixed by the lattice geometry alone, the relativistic factor
$\gamma = 1 + k\cdot\Delta\SSV$ is \emph{fully geometric} --- no free
parameter remains.  The Monte-Carlo verification routine in the GitHub
repository (Version~16, Appendix~A.8.1) recovers the exact low-velocity
expansion and saturation behavior with $k$ fixed at this value.

\subsection{Fractional quark charges as a cage-occupation prediction}
\label{sec:cage}

The elementary CP charge is $\pm 1$ and the lattice stiffness $C$ is
identical for every CP type.  The apparent fractional charges of quarks
($\pm 1/3$, $\pm 2/3$) therefore cannot arise from any difference in
fundamental coupling; they must be a projection effect at the composite level.

Standard Model particles in CPP are CP aggregates trapped in specific
600-cell symmetry cages (Version~16, \S 2 and Appendix~G):
\begin{itemize}
  \item \textbf{Lepton cage} (electron-type): a symmetry-protected orbit
    whose net external charge projects to $\pm 1$.
  \item \textbf{Quark cages} (up/down-type): lower-symmetry sub-cages whose
    internal CP orientations partially cancel under the H4 action, projecting
    to effective external charges of $\pm 1/3$ or $\pm 2/3$.
\end{itemize}
The specific fractions $1/3$ and $2/3$ are predicted (not inserted) by the
branching rules of the H4 Coxeter group acting on the 120-vertex
representation --- the same representation already used in Version~16
Appendix~G to select the 600-cell.  A full derivation of the branching table
and the resulting charge spectrum is reserved for the companion paper on
Standard Model particle charges; the present paper establishes only that the
\emph{mechanism} is cage-occupation geometry and that no new coupling
parameter is required.

\begin{remark}[Testable prediction]
If the cage-occupation prediction is correct, the ratio of quark to lepton
charge must be exactly $1/3$, and no intermediate values are possible for
color-singlet composites.  This is consistent with all experimental data.
Any observation of a non-$1/3$-quantized free charge would falsify this
sector of CPP.
\end{remark}

\subsection{Path to Newtonian gravity}

The same stiffness $C$ and 12-neighbor shell-broadcast mechanism that produce
the Coulomb force also produce Newtonian gravity when the SSV source is
mass-energy density rather than charge.  The gravitational constant $G$
emerges from the identical stiffness integral with the appropriate
Planck-scale normalization.  This is the subject of the next companion paper
in the CPP series.

% ══════════════════════════════════════════════════════════════════════════════
\section{Summary}
\label{sec:summary}
% ══════════════════════════════════════════════════════════════════════════════

We have established the microscopic origin of the Dipole Sea stiffness $C$:

\begin{enumerate}
  \item \textbf{Single effective medium.}  Equal abundances of all four DP
    types, identical elementary CP coupling strength, and quasi-static ZBW
    averaging reduce the four-component DP sea to a single effective stiffness
    (Section~\ref{sec:dp-sea}).

  \item \textbf{H4 Voronoi integral.}  That stiffness is the second-moment
    integral of the 600-cell Voronoi face-area distribution, evaluated exactly
    from the H4 golden-ratio coordinates in Version~16 Appendix~A.5, giving
    $C = \alphageom\cdot\EP/\lP^3$ with $\alphageom \approx 0.5594$
    (Section~\ref{sec:voronoi}).

  \item \textbf{Universal force law and charge spectrum.}  The same $C$ and
    shell-broadcast geometry reproduce the $1/r^2$ Coulomb law, fix the PSR
    formula without free parameters, and predict fractional quark charges as
    a cage-occupation effect of H4 representation theory
    (Section~\ref{sec:implications}).
\end{enumerate}

No new postulates are introduced.  The electromagnetic sector of CPP is now
closed: one geometric constant $\alphageom$, derived from the same 600-cell
that drives special relativity, fixes the universal charge coupling, the force
law, and the charge spectrum.

% ══════════════════════════════════════════════════════════════════════════════
\section*{References}
% ══════════════════════════════════════════════════════════════════════════════

\begin{enumerate}[label={[\arabic*]}]
  \item Abshier, T.~L. (2026).
    Mechanistic Derivation of Relativistic Effects via Space Stress Vector
    (SSV) in the Dipole Sea.  Version~16 (main paper).

  \item Abshier, T.~L. (2026).
    The Absolute Moment Postulate: Necessity, Consistency, and the PCD Cycle
    in Conscious Point Physics.  Version~3 (companion paper).

  \item Abshier, T.~L. (2025--2026).
    Conscious Point Physics Series.
    \url{https://hyperphysics.com/papers}
\end{enumerate}

\bigskip
\noindent\textbf{GitHub (Monte-Carlo stiffness verification):}\\
\url{https://github.com/tlabshier/CPP/blob/main/600-cell_special-relativity_emergence/600cell_monte_carlo_voronoi_k_fit.py}

\end{document}
